\section{Outcome (b): Combinatorics}

\subsection{What the outcome asks}

The exam tests classical probability: when every outcome of an experiment is equally likely,
\[
\Prob(A) = \frac{\text{number of outcomes in } A}{\text{total number of outcomes}},
\]
so the whole problem reduces to counting two sets correctly and consistently. The questions
ask you to decide whether order matters (permutation versus combination), whether selection is
with or without replacement, and how to split a group into several labelled subgroups
(multinomial partitions). The typical setting is selecting policies from a portfolio, forming
a committee from a group of policyholders, assigning claims to adjusters, or arranging files
in a queue. ``At least one'' questions are almost always fastest through the complement.
Exam arithmetic is done on a calculator, so answers are usually given to three decimals.

\subsection{Counting tools}

\paragraph{Multiplication principle.}
If a procedure consists of $k$ stages and stage $i$ can be done in $n_i$ ways regardless of
what happened at the earlier stages, the procedure can be done in $n_1 n_2 \cdots n_k$ ways.
Every formula below is a consequence of this principle.

\begin{keyfact}[Permutations: order matters, no replacement]
The number of ordered arrangements of all $n$ distinct objects is
\[
n! = n(n-1)(n-2)\cdots 2 \cdot 1, \qquad 0! = 1 .
\]
The number of ordered arrangements of $k$ objects chosen from $n$ distinct objects is
\[
{}_nP_k = \frac{n!}{(n-k)!} = n(n-1)\cdots(n-k+1) .
\]
\end{keyfact}

\begin{keyfact}[Combinations: order does not matter, no replacement]
The number of unordered subsets of size $k$ from $n$ distinct objects is
\[
\binom{n}{k} = \frac{n!}{k!\,(n-k)!} = \frac{{}_nP_k}{k!} .
\]
Useful identities:
\[
\binom{n}{k} = \binom{n}{n-k}, \qquad
\binom{n}{0} = \binom{n}{n} = 1, \qquad
\binom{n}{1} = n, \qquad
\binom{n}{2} = \frac{n(n-1)}{2},
\]
and Pascal's rule
\[
\binom{n}{k} = \binom{n-1}{k-1} + \binom{n-1}{k} .
\]
\end{keyfact}

Pascal's rule has a counting interpretation that recurs in exam problems: fix one particular
object; the $k$-subsets that contain it number $\binom{n-1}{k-1}$ and the $k$-subsets that
do not contain it number $\binom{n-1}{k}$. Dividing, the probability that a random $k$-subset
contains a specific object is $\binom{n-1}{k-1}/\binom{n}{k} = k/n$.

\begin{keyfact}[Binomial theorem]
\[
(x+y)^n = \sum_{k=0}^{n} \binom{n}{k} x^k y^{n-k}, \qquad
\sum_{k=0}^{n}\binom{n}{k} = 2^n, \qquad
\sum_{k=0}^{n}(-1)^k\binom{n}{k} = 0 .
\]
The identity $\sum_k \binom{n}{k}=2^n$ counts all subsets of an $n$-set. It is also why the
binomial probabilities in Topic 2 sum to one.
\end{keyfact}

\begin{keyfact}[Permutations with repeated objects (multinomial coefficient)]
The number of distinguishable arrangements of $n$ objects of which $n_1$ are identical of
type 1, $n_2$ identical of type 2, \dots, $n_r$ identical of type $r$
(with $n_1+\cdots+n_r=n$) is
\[
\binom{n}{n_1,\,n_2,\,\dots,\,n_r} = \frac{n!}{n_1!\,n_2!\cdots n_r!} .
\]
The same number counts the ways to partition $n$ distinct objects into $r$ labelled groups
of sizes $n_1,\dots,n_r$. It equals the product of successive combinations
$\binom{n}{n_1}\binom{n-n_1}{n_2}\cdots\binom{n_r}{n_r}$.
\end{keyfact}

For example, the word MISSISSIPPI has $11!/(1!\,4!\,4!\,2!) = 34650$ distinguishable
arrangements, and $12$ distinct claims can be divided among three adjusters, four claims each,
in $12!/(4!\,4!\,4!) = 34650$ ways.

\begin{keyfact}[With replacement]
Choosing $k$ times from $n$ objects with replacement, order recorded: $n^k$ outcomes.
Choosing $k$ times from $n$ objects with replacement, order ignored (a multiset): $\binom{n+k-1}{k}$ outcomes.
The second count is rarely needed on the exam; the first is needed constantly, in particular
as the denominator in ``at least one match'' problems.
\end{keyfact}

\begin{center}
\begin{tabular}{llp{3.4cm}p{5.2cm}}
\toprule
Order matters? & Replacement? & Count & Typical wording \\
\midrule
Yes & No  & ${}_nP_k = \dfrac{n!}{(n-k)!}$ & arrange, rank, sequence, fill distinct positions \\[6pt]
Yes & Yes & $n^k$ & each of $k$ items independently takes one of $n$ values \\[4pt]
No  & No  & $\dbinom{n}{k}$ & select, choose, committee, sample, hand \\[6pt]
No  & Yes & $\dbinom{n+k-1}{k}$ & multiset: only how many of each type is recorded \\
\bottomrule
\end{tabular}
\\[4pt]
Decision table for choosing $k$ from $n$ distinct objects.
\end{center}

\paragraph{Circular arrangements.}
$n$ distinct objects around a circle, where only relative position matters, can be arranged in
$(n-1)!$ ways: fix one object to remove the rotational symmetry and arrange the remaining $n-1$
in a row. If the circle can also be flipped over (a necklace), divide again by $2$.

\paragraph{Stars and bars.}
The number of ways to write $n$ as an ordered sum of $r$ nonnegative integers
$x_1 + x_2 + \cdots + x_r = n$, $x_i \ge 0$, is
\[
\binom{n+r-1}{r-1} .
\]
Picture $n$ identical stars and $r-1$ bars; each arrangement of the $n+r-1$ symbols is one
solution. If every $x_i \ge 1$ is required, give each variable one unit first and count
solutions of $y_1+\cdots+y_r = n-r$: $\binom{n-1}{r-1}$. For example, $10$ identical bonus
units can be distributed among $4$ agents in $\binom{13}{3}=286$ ways, and in
$\binom{9}{3}=84$ ways if every agent gets at least one. If all $286$ distributions are
equally likely, the probability that every agent gets at least one unit is $84/286 = 0.294$.

\paragraph{Restrictions.}
Two standard devices handle restrictions on arrangements.
\begin{itemize}[nosep]
\item \textbf{Objects that must be together}: glue them into one block, arrange the blocks,
then multiply by the number of internal orderings of the block. Three files that must be
adjacent among $7$ files in a row: $5!\cdot 3! = 720$ arrangements out of $7! = 5040$.
\item \textbf{Objects that must be separated}: arrange the unrestricted objects first, then
place the restricted objects in the gaps (including the two ends). No two of $3$ special
files adjacent among $7$: arrange the $4$ others ($4!$), choose $3$ of the $5$ gaps
($\binom{5}{3}$), order the special files in those gaps ($3!$): $24\cdot 10\cdot 6 = 1440$.
\end{itemize}

\subsection{Classical probability with counting}

With equally likely outcomes, $\Prob(A) = |A|/|S|$. The numerator and denominator must be
counted in the same way: both as ordered lists, or both as unordered sets. Either convention
gives the same probability if applied consistently.

\begin{keyfact}[The hypergeometric counting pattern]
A group contains $a$ objects of type I and $b$ objects of type II. A sample of $n$ objects is
drawn at random without replacement. The probability that the sample contains exactly $x$ of
type I is
\[
\Prob(X = x) = \frac{\dbinom{a}{x}\dbinom{b}{n-x}}{\dbinom{a+b}{n}} .
\]
With more than two types, the numerator becomes a product of one binomial coefficient per
type, and the denominator is $\binom{N}{n}$ with $N$ the total. Topic 2 names this the
hypergeometric distribution; on this outcome it is a counting pattern.
\end{keyfact}

\paragraph{Same problem, two methods.}
A portfolio holds $10$ policies: $4$ auto and $6$ homeowners. Three policies are selected at
random without replacement. Find the probability that exactly $2$ of the $3$ are auto.

\emph{Counting (unordered).} Total $\binom{10}{3} = 120$. Favourable
$\binom{4}{2}\binom{6}{1} = 6\cdot 6 = 36$. Probability $36/120 = 0.300$.

\emph{Sequential conditional probabilities (ordered).} The sequence auto, auto, home has
probability
\[
\frac{4}{10}\cdot\frac{3}{9}\cdot\frac{6}{8} = \frac{72}{720} = 0.100 .
\]
The sequences auto, home, auto and home, auto, auto each have the same probability (the
numerators $4,3,6$ appear in a different order, the denominators are $10,9,8$ regardless).
There are $\binom{3}{2}=3$ positions for the home policy, so the probability is
$3 \times 0.100 = 0.300$.

The two methods agree. The counting method is faster when the sample is unordered; the
sequential method is faster when the question involves only two or three draws, or asks
about a specific position (``the third policy selected is auto'').

\subsection{Worked examples}

\begin{example}[Committee with a constraint]
A group of $12$ policyholders consists of $7$ who hold auto coverage only and $5$ who hold
both auto and homeowners coverage. A committee of $5$ is chosen at random from the group.
Calculate the probability that the committee includes at least $2$ policyholders who hold
both coverages.
\par\smallskip\noindent (A) 0.44 \quad (B) 0.61 \quad (C) 0.69 \quad (D) 0.75 \quad (E) 0.81
\end{example}
\begin{solution}
Total committees: $\binom{12}{5} = 792$. Use the complement: fewer than $2$ from the
``both'' group means $0$ or $1$.
\[
\binom{5}{0}\binom{7}{5} + \binom{5}{1}\binom{7}{4} = 21 + 5\cdot 35 = 196 .
\]
\[
\Prob(\text{at least } 2) = 1 - \frac{196}{792} = \frac{596}{792} = 0.7525 .
\]
Answer: \textbf{(D)}.

The direct sum $\binom{5}{2}\binom{7}{3}+\binom{5}{3}\binom{7}{2}+\binom{5}{4}\binom{7}{1}+\binom{5}{5}\binom{7}{0}
= 350+210+35+1 = 596$ agrees. Choice (A) is the probability of exactly $2$, choice (C) is
the probability of at most $2$.
\end{solution}

\begin{example}[Sampling a portfolio with three policy types]
An insurer's portfolio consists of $20$ policies: $8$ auto, $7$ homeowners, and $5$ life.
Four policies are selected at random without replacement for audit. Calculate the probability
that the audit sample contains exactly $2$ auto, $1$ homeowners, and $1$ life policy.
\par\smallskip\noindent (A) 0.006 \quad (B) 0.058 \quad (C) 0.202 \quad (D) 0.263 \quad (E) 0.350
\end{example}
\begin{solution}
\[
\Prob = \frac{\binom{8}{2}\binom{7}{1}\binom{5}{1}}{\binom{20}{4}}
= \frac{28\cdot 7\cdot 5}{4845} = \frac{980}{4845} = 0.2023 .
\]
Answer: \textbf{(C)}.

Choice (B) results from writing $8\cdot 7\cdot 5$ in the numerator (forgetting that the
two auto policies are an unordered pair, $\binom{8}{2}=28$, not $8\cdot 7$ or $8$).
Choice (A) results from an ordered numerator $8\cdot7\cdot5\cdot3$ over the ordered
denominator ${}_{20}P_4$ without multiplying by the $4!/2! = 12$ arrangements of the types.
\end{solution}

\begin{example}[At least one match: birthday pattern]
A claims department has $8$ adjusters. Each of $4$ incoming claims is assigned at random to
one of the $8$ adjusters, independently of the other claims. Calculate the probability that at
least two of the four claims are assigned to the same adjuster.
\par\smallskip\noindent (A) 0.41 \quad (B) 0.50 \quad (C) 0.55 \quad (D) 0.59 \quad (E) 0.66
\end{example}
\begin{solution}
Assignments are with replacement and ordered (claim 1 to adjuster $x$, claim 2 to adjuster
$y$, \dots), so the sample space has $8^4 = 4096$ equally likely outcomes. The complement of
``at least two share'' is ``all four different'', with ${}_8P_4 = 8\cdot7\cdot6\cdot5 = 1680$
outcomes.
\[
\Prob(\text{at least one shared adjuster}) = 1 - \frac{1680}{4096} = 1 - 0.4102 = 0.5898 .
\]
Answer: \textbf{(D)}.

The same structure gives the birthday problem: with $23$ people and $365$ equally likely
birthdays, $1 - {}_{365}P_{23}/365^{23} = 0.507$.
\end{solution}

\begin{example}[Arrangement with a separation restriction]
Seven claim files are placed in a queue in random order. Three of the files are auto claims
and four are homeowners claims. Calculate the probability that no two auto claim files are
adjacent in the queue.
\par\smallskip\noindent (A) $1/7$ \quad (B) $2/7$ \quad (C) $3/7$ \quad (D) $4/7$ \quad (E) $5/7$
\end{example}
\begin{solution}
Treat the files as distinct. Total orderings: $7! = 5040$. Arrange the four homeowners files
first: $4! = 24$ ways. They create $5$ gaps (three interior, two ends). Choose $3$ of the $5$
gaps for the auto files, $\binom{5}{3} = 10$, and order the auto files within the chosen gaps,
$3! = 6$. Favourable orderings: $24\cdot 10\cdot 6 = 1440$.
\[
\Prob = \frac{1440}{5040} = \frac{2}{7} .
\]
Answer: \textbf{(B)}.

The same answer is reached treating the files as indistinguishable within type: the queue is
a pattern of $3$ A's and $4$ H's, $\binom{7}{3} = 35$ patterns, of which $\binom{5}{3} = 10$
have no two A's adjacent; $10/35 = 2/7$. Either model works if used for both numerator and
denominator. For contrast, the probability that the three auto files are all adjacent is
$5!\cdot 3!/7! = 720/5040 = 1/7$, choice (A).
\end{solution}

\begin{example}[Multinomial partition of claims among adjusters]
Twelve claims of different sizes are to be divided at random among three adjusters, with each
adjuster receiving exactly four claims. Calculate the probability that the three largest
claims are all assigned to the same adjuster.
\par\smallskip\noindent (A) 0.018 \quad (B) 0.027 \quad (C) 0.036 \quad (D) 0.045 \quad (E) 0.055
\end{example}
\begin{solution}
\emph{Counting.} Total partitions: $\dfrac{12!}{4!\,4!\,4!} = 34650$. Favourable: choose which
adjuster receives the three largest claims ($3$ ways), choose that adjuster's fourth claim from
the remaining $9$ ($\binom{9}{1}=9$), then split the remaining $8$ claims between the other two
adjusters ($\binom{8}{4} = 70$). Favourable count $3\cdot 9\cdot 70 = 1890$.
\[
\Prob = \frac{1890}{34650} = \frac{3}{55} = 0.0545 .
\]
\emph{Sequential.} Wherever the largest claim goes, that adjuster has $3$ open slots among
the $11$ remaining; the second largest lands there with probability $3/11$, and then the
third largest with probability $2/10$: $\dfrac{3}{11}\cdot\dfrac{2}{10} = \dfrac{6}{110} = \dfrac{3}{55}$.

Answer: \textbf{(E)}.

Choice (A) is $\binom{4}{3}/\binom{12}{3}$, which is the probability that a specified adjuster
gets all three; multiplying by $3$ adjusters gives the correct $0.055$.
\end{solution}

\begin{example}[Solving for the group size]
An insurer will audit $2$ of its $n$ agents, chosen at random. There are exactly $45$
different pairs of agents that could be audited. The insurer instead decides to audit $3$
agents chosen at random. Calculate the probability that a particular agent is among the three
audited.
\par\smallskip\noindent (A) 0.10 \quad (B) 0.20 \quad (C) 0.30 \quad (D) 0.33 \quad (E) 0.45
\end{example}
\begin{solution}
$\binom{n}{2} = \dfrac{n(n-1)}{2} = 45$, so $n(n-1) = 90 = 10\cdot 9$ and $n = 10$
(the other root of $n^2 - n - 90 = 0$ is $n=-9$, rejected).

With $n=10$, the number of $3$-subsets containing the particular agent is
$\binom{9}{2} = 36$ (the agent is in; choose $2$ companions from the other $9$), out of
$\binom{10}{3} = 120$. Probability $36/120 = 0.30$. Equivalently, $k/n = 3/10$.

Answer: \textbf{(C)}.
\end{solution}

\begin{example}[With replacement: every option chosen at least once]
A survey asks $5$ customers to select one of $3$ deductible options. Each customer selects an
option at random, independently, with each option equally likely. Calculate the probability
that every option is selected by at least one customer.
\par\smallskip\noindent (A) 0.383 \quad (B) 0.617 \quad (C) 0.667 \quad (D) 0.753 \quad (E) 0.802
\end{example}
\begin{solution}
Ordered outcomes with replacement: $3^5 = 243$. Let $A_i$ be the event that option $i$ is
selected by nobody. By inclusion--exclusion,
\[
\Prob(A_1\cup A_2\cup A_3) = 3\left(\frac{2}{3}\right)^5 - 3\left(\frac{1}{3}\right)^5 + 0
= \frac{3\cdot 32 - 3\cdot 1}{243} = \frac{93}{243} .
\]
\[
\Prob(\text{every option used}) = 1 - \frac{93}{243} = \frac{150}{243} = 0.6173 .
\]
Answer: \textbf{(B)}.

Direct count: the $5$ customers split over the $3$ options as $3{+}1{+}1$ or $2{+}2{+}1$.
Pattern $3{+}1{+}1$: choose the option used three times ($3$), choose those three customers
($\binom{5}{3} = 10$), assign the remaining two customers to the two remaining options
($2! = 2$): $60$. Pattern $2{+}2{+}1$: choose the option used once ($3$), choose that customer
($5$), split the remaining four customers into two labelled pairs ($\binom{4}{2} = 6$): $90$.
Total $150$, as before. Choice (A) is the probability of the complement.
\end{solution}

\subsection{Traps}

\begin{trap}[Ordered numerator over unordered denominator, or the reverse]
The most common combinatorics error is counting the favourable outcomes as ordered sequences
(products such as $8\cdot 7\cdot 5$) while the denominator is an unordered count
$\binom{N}{n}$, or the reverse. Decide at the start whether the sample space is ordered or
unordered and count both sets the same way. When in doubt, use unordered counts with
$\binom{\cdot}{\cdot}$ throughout.
\end{trap}

\begin{trap}[Double counting in ``at least $k$'' problems]
To count committees with at least $2$ from a group of $5$, it is wrong to ``choose $2$ from
the $5$, then choose the other $3$ from the remaining $10$'': $\binom{5}{2}\binom{10}{3} = 1200$
exceeds the total $\binom{12}{5}=792$, because a committee with three from the group is counted
three times. Either sum the exact cases ($\binom{5}{2}\binom{7}{3} + \binom{5}{3}\binom{7}{2} + \cdots$)
or use the complement.
\end{trap}

\begin{trap}[Forgetting to divide by $k!$]
${}_nP_k$ counts ordered selections; $\binom{n}{k} = {}_nP_k/k!$ counts unordered ones. A
question that asks for a committee, a sample, or a hand wants $\binom{n}{k}$. Writing
$8\cdot 7 = 56$ for the number of pairs of auto policies, when the number is
$\binom{8}{2} = 28$, doubles the numerator.
\end{trap}

\begin{trap}[Identical objects treated as distinct, or the reverse]
The word ``identical'' or ``indistinguishable'' in a question signals a multinomial
coefficient $n!/(n_1!\cdots n_r!)$, not $n!$. Conversely, when policies or claims are described
by type but are physically distinct items, they may be treated as distinct, provided the
denominator uses the same convention. The safe approach for probability questions is to treat
all objects as distinct in both numerator and denominator.
\end{trap}

\begin{trap}[``At least one'' by direct summation]
Summing the cases $1, 2, \dots, n$ for ``at least one'' wastes time and multiplies the chances
of an arithmetic slip. Compute $1 - \Prob(\text{none})$. For ``at least two'', compute
$1 - \Prob(0) - \Prob(1)$.
\end{trap}

\subsection{Practice problems}

\begin{practice}
A claims department has $8$ senior adjusters and $6$ junior adjusters. A review team of $4$ is
selected at random from the $14$. Calculate the probability that the team consists of exactly
$2$ senior and $2$ junior adjusters.
\par\smallskip\noindent (A) 0.21 \quad (B) 0.32 \quad (C) 0.42 \quad (D) 0.51 \quad (E) 0.60
\end{practice}

\begin{practice}
Of $20$ claims submitted in a month, $3$ are fraudulent. An auditor selects $5$ of the $20$
claims at random without replacement. Calculate the probability that at least one fraudulent
claim is selected.
\par\smallskip\noindent (A) 0.40 \quad (B) 0.52 \quad (C) 0.56 \quad (D) 0.60 \quad (E) 0.75
\end{practice}

\begin{practice}
Five auto claim files and three homeowners claim files are placed in a processing queue in
random order. Calculate the probability that the three homeowners files are processed
consecutively.
\par\smallskip\noindent (A) $1/56$ \quad (B) $3/28$ \quad (C) $1/8$ \quad (D) $3/14$ \quad (E) $3/8$
\end{practice}

\begin{practice}
The number of distinct pairs of policies that can be formed from a portfolio of $n$ policies
is $66$. Four policies are selected at random from the portfolio. Calculate the probability
that two specified policies are both among the four selected.
\par\smallskip\noindent (A) $1/66$ \quad (B) $1/33$ \quad (C) $1/22$ \quad (D) $1/12$ \quad (E) $1/11$
\end{practice}

\begin{practice}
Nine claims are divided at random among three adjusters so that each adjuster receives
exactly three claims. Calculate the probability that two particular claims are assigned to
the same adjuster.
\par\smallskip\noindent (A) $1/4$ \quad (B) $2/7$ \quad (C) $1/3$ \quad (D) $4/9$ \quad (E) $1/2$
\end{practice}

\begin{practice}
Four customers each independently select one of $5$ deductible options, with every option
equally likely. Calculate the probability that the four customers select four different
options.
\par\smallskip\noindent (A) 0.096 \quad (B) 0.192 \quad (C) 0.288 \quad (D) 0.384 \quad (E) 0.480
\end{practice}

\begin{practice}
A pricing committee of $4$ is selected at random from $5$ actuaries and $4$ underwriters.
Calculate the probability that the committee includes at least one actuary and at least one
underwriter.
\par\smallskip\noindent (A) 0.833 \quad (B) 0.881 \quad (C) 0.905 \quad (D) 0.952 \quad (E) 0.976
\end{practice}

\subsection{Answers to practice problems}

\begin{enumerate}[label=\arabic*.,nosep]
\item \textbf{(C)}. $\binom{8}{2}\binom{6}{2}/\binom{14}{4} = 28\cdot 15/1001 = 420/1001 = 0.4196$.
\item \textbf{(D)}. $1 - \binom{17}{5}/\binom{20}{5} = 1 - 6188/15504 = 1 - 0.3991 = 0.6009$.
Choice (C), $1-(17/20)^5 = 0.556$, wrongly assumes replacement.
\item \textbf{(B)}. Glue the three homeowners files into one block: $6!$ arrangements of six
blocks, times $3!$ internal orders, over $8!$: $6\cdot 720/40320 = 4320/40320 = 3/28$.
\item \textbf{(E)}. $\binom{n}{2}=66 \Rightarrow n(n-1)=132=12\cdot 11$, $n=12$. Both specified
policies in: $\binom{10}{2}/\binom{12}{4} = 45/495 = 1/11$. Sequentially:
$\frac{4}{12}\cdot\frac{3}{11} = 1/11$.
\item \textbf{(A)}. Wherever the first claim goes, that adjuster has $2$ open slots among the
$8$ remaining claims: $2/8 = 1/4$. Counting: $3\cdot\binom{7}{1}\binom{6}{3}/\dfrac{9!}{3!\,3!\,3!} = 420/1680 = 1/4$.
\item \textbf{(B)}. ${}_5P_4/5^4 = 120/625 = 0.192$.
\item \textbf{(D)}. Complement: all actuaries or all underwriters,
$\left[\binom{5}{4}+\binom{4}{4}\right]/\binom{9}{4} = 6/126$; probability $120/126 = 20/21 = 0.952$.
\end{enumerate}

\subsection{One-page summary}

\begin{itemize}[nosep]
\item Classical probability: $\Prob(A) = |A|/|S|$ for equally likely outcomes. Count numerator
and denominator by the same convention (both ordered or both unordered).
\item Multiplication principle: $n_1 n_2\cdots n_k$ for a $k$-stage procedure.
\item Ordered, no replacement: ${}_nP_k = n!/(n-k)!$; all $n$: $n!$.
\item Unordered, no replacement: $\binom{n}{k} = n!/[k!(n-k)!] = {}_nP_k/k!$.
\item Ordered, with replacement: $n^k$. Unordered, with replacement: $\binom{n+k-1}{k}$.
\item Identities: $\binom{n}{k}=\binom{n}{n-k}$; $\binom{n}{2}=n(n-1)/2$;
$\binom{n}{k}=\binom{n-1}{k-1}+\binom{n-1}{k}$; $\sum_k\binom{n}{k}=2^n$;
$(x+y)^n=\sum_k\binom{n}{k}x^ky^{n-k}$.
\item A random $k$-subset of $n$ contains a specified object with probability $k/n$.
\item Multinomial coefficient $n!/(n_1!\cdots n_r!)$: arrangements with identical objects, and
partitions of $n$ distinct objects into labelled groups of sizes $n_1,\dots,n_r$.
\item Hypergeometric pattern: $\binom{a}{x}\binom{b}{n-x}/\binom{a+b}{n}$; extend with one
factor per type. Equals the sequential product of conditional probabilities times the number
of orderings.
\item Circular arrangements: $(n-1)!$. Stars and bars: $x_1+\cdots+x_r=n$, $x_i\ge0$, has
$\binom{n+r-1}{r-1}$ solutions; $x_i\ge1$ has $\binom{n-1}{r-1}$.
\item Must be adjacent: glue into a block, arrange blocks, multiply by internal orderings.
Must be separated: arrange the others, then choose gaps.
\item ``At least one'': $1-\Prob(\text{none})$. Birthday pattern: $1-{}_nP_k/n^k$.
\item Solve $\binom{n}{2}=m$ by $n(n-1)=2m$ and factoring into consecutive integers.
\item Never ``fix $k$ then choose the rest'' for ``at least $k$''; sum exact cases or use
the complement.
\end{itemize}
